\documentclass[11pt,a4paper]{article}
\usepackage[T1]{fontenc}
\usepackage{lmodern}
\usepackage[margin=25mm]{geometry}
\usepackage{amsmath,amssymb,amsthm,mathtools,mathrsfs}
\usepackage{booktabs,longtable,array,microtype,xurl}
\usepackage[colorlinks=true,linkcolor=blue!45!black,citecolor=blue!45!black,urlcolor=blue!45!black]{hyperref}
\usepackage{xcolor,fancyhdr,bookmark}
\setlength{\headheight}{14pt}
\pagestyle{fancy}\fancyhf{}
\fancyhead[L]{\small Discrete vertical Hurwitz calculus}
\fancyhead[R]{\small ProveIt continuation}
\fancyfoot[C]{\thepage}
\setlength{\emergencystretch}{3em}
\setlength{\parskip}{3pt}
\allowdisplaybreaks[2]
\numberwithin{equation}{section}
\newtheorem{theorem}{Theorem}[section]
\newtheorem{lemma}[theorem]{Lemma}
\newtheorem{proposition}[theorem]{Proposition}
\newtheorem{corollary}[theorem]{Corollary}
\theoremstyle{definition}
\newtheorem{definition}[theorem]{Definition}
\theoremstyle{remark}
\newtheorem{remark}[theorem]{Remark}
\newtheorem{example}[theorem]{Example}
\DeclareMathOperator{\Li}{Li}
\DeclareMathOperator{\Log}{Log}
\DeclareMathOperator{\Ree}{Re}
\DeclareMathOperator{\CT}{CT}
\DeclareMathOperator{\Sym}{Sym}
\newcommand{\dd}{\,\mathrm d}
\newcommand{\ii}{\mathrm i}
\newcommand{\Z}{\mathbb Z}
\newcommand{\N}{\mathbb N}
\newcommand{\C}{\mathbb C}
\newcommand{\eps}{\varepsilon}
\newcommand{\Rcal}{\mathcal R}
\newcommand{\Bcal}{\mathcal B}
\newcommand{\Lcal}{\mathcal L}
\newcommand{\Qcal}{\mathcal Q}
\newcommand{\Scal}{\mathcal S}
\newcommand{\Ecal}{\mathcal E}
\newcommand{\src}[1]{\nolinkurl{#1}}
\hypersetup{pdftitle={Discrete Vertical Hurwitz Calculus: Stieltjes Jets, Polylogarithmic Sampling, and Resonant Contact Terms},pdfauthor={Research continuation for the ProveIt project}}
\title{\textbf{Discrete Vertical Hurwitz Calculus}\\[7pt]
\large Stieltjes jets, polylogarithmic sampling,\\and resonant contact terms}
\author{Research continuation prepared for\\Vladimir Reshetnikov's ProveIt project}
\date{11 October 2026}
\begin{document}
\maketitle
\begin{abstract}
We develop an exact calculus for sums over a vertical lattice in the
Hurwitz argument. Euler--Maclaurin compensation and an elementary
periodization argument give a normally convergent spectral identity
whose right-hand side is a Lambert series of polylogarithms. Its complete
Taylor expansion at order one evaluates every compensated generalized
Stieltjes lattice sum, with finite Bell-polynomial coefficients. Horizontal
unit differences eliminate the Lambert series and leave finitely many
polylogarithms and their order derivatives. Anchored antiderivatives are
obtained by raising the polylogarithm order.

At nonpositive integer orders we prove a parity-sensitive contact law.
First spectral derivatives have absolutely convergent lattice sums at
negative odd orders; at negative even orders the symmetric sum differs
from the spectral continuation by an explicit Bernoulli constant. The
Gamma case is an already published specialization of Arabi Ardehali and
Rosengren, not a novelty claim. All higher, sufficiently compensated
spectral jets are expressed by a finite compiler involving generalized
harmonic numbers and zeta values. Mean-zero periodic lattice weights
annihilate every first-jet contact.

A second theorem evaluates lattice products of arbitrary positive-order
polygammas as a single Hurwitz term plus exponentially convergent image
sums, each image being a finite Hurwitz--Lerch combination. We also give
an all-index Stieltjes-difference pairing and an explicit four-shift
digamma formula. Proofs distinguish ordinary absolute sums, symmetric
limits, and analytic continuation. The package includes 268 exact
algebraic assertions, 38 independent or cross-route numerical comparisons,
two corruption controls, and a scoped source audit. The numerical tests
are not interval certificates. No claim is made to settle the manuscript's
$S_6$, revised $S_8$, or general unequal-tail prefix conjectures.
\end{abstract}
\clearpage
\tableofcontents
\clearpage

\section{Scope, source frontier, and contribution}

The ProveIt manuscript has developed several complementary calculi:
periodic Hurwitz correlations, spectral derivatives, finite parts,
Gamma products, and harmonic-tail reconstruction. In particular, the
inspected vertical-pairing report concerns ordinary integrals in the
imaginary part of the Hurwitz argument, rather than discrete sampling
\cite{RepoVertical}. The present article replaces that continuous
vertical geometry by the lattice
\[
 a+\ii h\Z,\qquad h>0,\qquad \Ree a>0.
\]
This change produces nonzero sampling images as well as an explicit
boundary contact. It therefore requires its own convergence statements
and cannot be obtained by simply replacing an integral sign by a sum.

The principal exact results are the compensated sampling theorem
(Theorem~\ref{thm:sampling}), the complete Stieltjes coefficient identity
(Theorem~\ref{thm:stieltjes}), the negative-integer contact law
(Theorem~\ref{thm:negative}), the harmonic-number jet compiler
(Theorem~\ref{thm:negative-jets}), the finite-image polygamma pairing
(Theorem~\ref{thm:bilinear}), and the four-shift formula
(Theorem~\ref{thm:fourshift}). The proofs are independent of the pending
harmonic-rank conjectures.

\paragraph{What is inherited.}
The Hurwitz continuation, Euler--Maclaurin expansion, Gamma special
values, and polylogarithm/Lerch definitions are classical
\cite{DLMF25,DLMF5,DLMFPoly,DLMFLerch}. Periodization is a standard
form of Poisson summation \cite{DLMFFourier}. Stieltjes summation and
parameter differentiation also have substantial antecedents; Coffey's
summation relations are one relevant example \cite{Coffey}. Most
importantly, the Gamma product obtained below is equivalent to the
published product formula of Arabi Ardehali and Rosengren
\cite[Eq.~(13)]{ArdehaliRosengren}. That specialization, including its
constant, is explicitly credited.

\paragraph{What is proposed for integration.}
The contribution is the unified discrete spectral calculus, its
all-index convergent Stieltjes and higher-resonance consequences, and the
explicit image formulas for products. These are proved here as exact
identities. They are proposed continuations relative to the inspected
repository sources, not assertions of exhaustive historical priority.
The distinction matters: a preliminary reconstruction calculation in
this investigation reproduced the already proved smaller-tail-order-two
family in a later repository report \cite{RepoBernoulli}. It is not
included as a new result.

\paragraph{Audit limits.}
The main branch changed during inspection. Accordingly, the audit
records individual source paths and blob identifiers rather than
claiming a single immutable snapshot for all reads. The incoming
directory was listed and its intake instructions were read, but the
interiors of its five then-listed ZIP archives were not available for
inspection in this run. No complete audit of those archives, of all
manuscript volumes, or of the literature is claimed. No repository file
was modified. Section~\ref{sec:audit} records the concrete correction
and the distinction between that correction and a mathematical warning.

\section{Conventions and the compensated Mellin family}
\label{sec:conventions}

Fix $h>0$ and write
\begin{equation}\label{eq:scale}
 \tau=\frac{2\pi}{h},\qquad q=e^{-\tau},\qquad z=e^{-\tau a},
 \qquad w_n=a+\ii hn\quad(n\in\Z).
\end{equation}
Initially $\Ree a>0$; hence $|z|<1$. Every power of $w$ uses the analytic
logarithm $\Log w$ in the right half-plane. The analytic branch of
$\log\Gamma(w)$ is real on the positive real axis. It is not the
principal logarithm of a numerically computed complex Gamma value.

We use the Laurent convention
\begin{equation}\label{eq:stieltjes-convention}
 \zeta(1+\eps,w)=\frac1\eps+
 \sum_{m=0}^{\infty}\frac{(-1)^m}{m!}\gamma_m(w)\eps^m,
 \qquad \gamma_0(w)=-\psi(w).
\end{equation}
Here $\psi=\Gamma'/\Gamma$, $(s)_k=s(s+1)\cdots(s+k-1)$,
$B_{2j}$ are Bernoulli numbers, and
\[
 H_N^{(r)}=\sum_{k=1}^N k^{-r},\qquad H_N=H_N^{(1)},\qquad H_0^{(r)}=0.
\]
The Bernoulli convention is $B_1=-1/2$; no formula below depends on a
software package's alternative value for $B_1$.

\begin{definition}[Compensation]\label{def:compensation}
For an integer $M\ge0$, put
\begin{align}
 C_M(t)&=\frac1{1-e^{-t}}-\frac1t-\frac12
       -\sum_{j=1}^M\frac{B_{2j}}{(2j)!}t^{2j-1},\label{eq:kernel}\\
 \Rcal_M(s,w)&=\zeta(s,w)-\frac{w^{1-s}}{s-1}-\frac12w^{-s}
       -\sum_{j=1}^M\frac{B_{2j}}{(2j)!}(s)_{2j-1}w^{1-s-2j}.
 \label{eq:remainder}
\end{align}
The singularity at $s=1$ in the combined expression is filled by its
removable limit. Thus $\Rcal_M$ is entire in $s$ for each $\Ree w>0$.
\end{definition}

The subtractions are fixed, explicit functions. No unspecified
regularization or finite-part convention is hidden in
\eqref{eq:remainder}. Near zero,
\begin{equation}\label{eq:kernel-origin}
 C_M(t)=\frac{B_{2M+2}}{(2M+2)!}t^{2M+1}+O(t^{2M+3}).
\end{equation}
For $\Ree s>-2M-1$, the Mellin identity is
\begin{equation}\label{eq:mellin}
 \Rcal_M(s,w)=\frac1{\Gamma(s)}
       \int_0^\infty t^{s-1}e^{-wt}C_M(t)\dd t.
\end{equation}
It follows initially from the ordinary Hurwitz integral in a common
convergence domain, by integrating the finitely many subtractions.
Both sides are holomorphic on the indicated half-plane after removable
values are filled, so the identity continues there. The factor
$1/\Gamma(s)$ is entire; at its zeros the convergent integral is
multiplied by zero.

\begin{lemma}[Uniform vertical remainder]\label{lem:vertical}
For $s$ in a compact set and $w$ in a fixed sector
$|\arg w|\le\pi-\delta$, the Euler--Maclaurin expansion gives, for
sufficiently large $J$,
\begin{equation}\label{eq:em-asymptotic}
 \Rcal_M(s,w)=
 \sum_{j=M+1}^{J}\frac{B_{2j}}{(2j)!}(s)_{2j-1}w^{1-s-2j}
 +O\bigl(|w|^{-\Ree s-2J-1}\bigr).
\end{equation}
In particular, on compact subsets of
$\{\Ree s>-2M,\ \Ree a>0\}$ the series
\[
 \sum_{n\in\Z}\partial_s^d\Rcal_M(s,a+\ii hn)
\]
converges normally for each fixed $d\ge0$.
\end{lemma}
\begin{proof}
Apply Euler--Maclaurin with additional terms to the defining Hurwitz
sum; its remainder is an integral of a bounded periodic Bernoulli
function times a power of $w+x$. In a closed sector avoiding the
negative axis, this gives the stated uniform expansion. Equivalently,
this is the standard sectorial Hurwitz expansion
\cite[\S25.11(xii)]{DLMF25}, with the first $M$ terms removed.
On a vertical line the first remaining term has size
$O(|n|^{-\Ree s-2M-1})$. Spectral derivatives add only fixed powers
of $\log(2+|n|)$, by Cauchy estimates on a slightly larger compact
set. The exponent exceeds one uniformly on the indicated compacts.
The finitely many small $n$ terms are holomorphic and bounded there.
\end{proof}

Define the entire-in-order Lambert--polylogarithm family
\begin{equation}\label{eq:Lambert}
 \Lcal_\nu(q,z)=\sum_{r=0}^{\infty}\Li_\nu(zq^r)
              =\sum_{k=1}^{\infty}\frac{z^k}{k^\nu(1-q^k)}.
\end{equation}
Both expressions converge normally for $|z|<1$ and $\nu$ in compact
subsets of $\C$. Every order derivative is obtained by inserting
$(-\log k)^d$ in the second series. We will use
\begin{align}
 \Bcal_M(s;a,\tau)={}&\Lcal_{1-s}(q,z)-\frac12\Li_{1-s}(z)
          -\frac1\tau\Li_{2-s}(z)\notag\\
 &-\sum_{j=1}^M\frac{B_{2j}\tau^{2j-1}}{(2j)!}\Li_{2-s-2j}(z).
 \label{eq:bracket}
\end{align}
It is entire in $s$. All $\tau$-dependence, including $q$ and $z$,
is understood as in \eqref{eq:scale}.

\section{An exact vertical sampling theorem}

\begin{lemma}[Periodization with absolutely summable coefficients]
\label{lem:periodization}
Let $f\in L^1(\mathbb R)$ be continuous and suppose
$P(x)=\sum_{k\in\Z}f(x+\tau k)$ converges locally uniformly to a
continuous $\tau$-periodic function. If
$\sum_n|\widehat f(hn)|<\infty$, with
\[
 \widehat f(y)=\int_{\mathbb R}f(t)e^{-\ii yt}\dd t,
 \qquad h\tau=2\pi,
\]
then
\begin{equation}\label{eq:periodization}
 \sum_{n\in\Z}\widehat f(hn)=\tau\sum_{k\in\Z}f(\tau k).
\end{equation}
\end{lemma}
\begin{proof}
The $n$th Fourier coefficient of $P$ is $\widehat f(hn)/\tau$,
by integration over one period and absolute integrability. Their
absolute summability gives a uniformly convergent Fourier series.
It represents $P$: the difference is a continuous periodic function
with every Fourier coefficient zero, hence is zero, for example by
Fej\'er approximation. Evaluate that Fourier series at $x=0$.
\end{proof}

\begin{theorem}[Compensated vertical-lattice sampling]
\label{thm:sampling}
Let $M\ge0$, $h>0$, $\Ree a>0$, and $\Ree s>-2M$. Then
\begin{equation}\label{eq:master}
 \boxed{\displaystyle
 \Scal_M(s;a,h):=\sum_{n\in\Z}\Rcal_M(s,a+\ii hn)
       =\frac{\tau^s}{\Gamma(s)}\Bcal_M(s;a,\tau).}
\end{equation}
The lattice sum is absolutely and normally convergent on this domain,
including every fixed spectral derivative. The right side gives an
entire continuation in $s$ of this summed function.
\end{theorem}
\begin{proof}
Extend
\[
 f(t)=\frac{t^{s-1}e^{-at}C_M(t)}{\Gamma(s)}\quad(t>0)
\]
by zero for $t\le0$. At zero it is continuous and vanishes, because
its leading power is $t^{s+2M}$ and $\Ree(s+2M)>0$.
At positive infinity it decays exponentially times a fixed power,
so its periodization is continuous and locally uniformly convergent.
Its Fourier transform is $\Rcal_M(s,a+\ii y)$ by
\eqref{eq:mellin}, and the Fourier coefficients are absolutely
summable by Lemma~\ref{lem:vertical}. Thus
\[
 \Scal_M(s;a,h)=\frac{\tau^s}{\Gamma(s)}
       \sum_{k=1}^\infty k^{s-1}z^k C_M(\tau k).
\]
Insert \eqref{eq:kernel} and identify each of the absolutely
convergent series. This gives \eqref{eq:bracket} and
\eqref{eq:master}. The normal convergence statement is
Lemma~\ref{lem:vertical}; the continuation follows from the
entire order dependence on the right.
\end{proof}

\begin{remark}[Why the strict boundary matters]\label{rem:boundary}
At $\Ree s=-2M$, the density in the proof need not vanish at zero.
The original lattice sum need not be absolutely convergent in a
neighborhood of that order. Consequently, differentiating the entire
right side there does not automatically differentiate an ordinary
lattice sum. Section~\ref{sec:negative} computes the resulting contact
exactly. This is a failure of interchange, not a failure of analytic
continuation.
\end{remark}

The theorem also covers an offset lattice $a+\ii h(n+\theta)$:
replace $a$ by $a+\ii h\theta$, so every polylogarithm argument is
multiplied by $e^{-2\pi\ii\theta}$. No new convergence issue arises.

\section{Every generalized Stieltjes lattice sum}
\label{sec:stieltjes}

Set
\begin{equation}\label{eq:K}
 K_m(w)=\gamma_m(w)+\frac{\Log^{m+1}w}{m+1}
                     -\frac{\Log^m w}{2w}.
\end{equation}
The exact Laurent cancellation gives
\begin{equation}\label{eq:K-generator}
 \Rcal_0(1+\eps,w)=\sum_{m=0}^\infty\frac{(-1)^m}{m!}K_m(w)\eps^m.
\end{equation}
In particular, each $K_m(a+\ii hn)$ is
$O(\log^m(2+|n|)/(1+|n|)^2)$.

For a variable $Y$, define finite Bell polynomials $\Ecal_j(Y)$ by
\begin{equation}\label{eq:E}
 \exp\left(Yu+\sum_{r=2}^\infty
             \frac{(-1)^{r+1}\zeta(r)}r u^r\right)
      =\sum_{j=0}^\infty\Ecal_j(Y)\frac{u^j}{j!}.
\end{equation}
Only $r\le j$ contributes to $\Ecal_j$. Their first values are
\begin{align*}
 \Ecal_0(Y)&=1,& \Ecal_1(Y)&=Y,\\
 \Ecal_2(Y)&=Y^2-\zeta(2),&
 \Ecal_3(Y)&=Y^3-3Y\zeta(2)+2\zeta(3),\\
 \Ecal_4(Y)&=Y^4-6Y^2\zeta(2)+8Y\zeta(3)
             +3\zeta(2)^2-6\zeta(4).
\end{align*}
Write
\begin{equation}\label{eq:Q}
 \Qcal_\nu(a)=\Lcal_\nu(q,z)-\frac12\Li_\nu(z)-\frac1\tau\Li_{\nu+1}(z),
 \qquad L=\log\tau+\gamma.
\end{equation}
Here $\gamma$ is Euler's constant, and $\tau$ is held fixed in every
$a$- or order derivative.

\begin{theorem}[All-index Stieltjes identity]\label{thm:stieltjes}
For $m\ge0$ and $\Ree a>0$,
\begin{equation}\label{eq:stieltjes-main}
 \boxed{\displaystyle
 \sum_{n\in\Z}K_m(a+\ii hn)
 =\tau\sum_{j=0}^m\binom mj(-1)^j\Ecal_j(L)
        \left.\partial_\nu^{\,m-j}\Qcal_\nu(a)\right|_{\nu=0}.}
\end{equation}
The sum is ordinary and absolutely convergent. There is no principal
value and no finite part.
\end{theorem}
\begin{proof}
In \eqref{eq:master}, put $M=0$ and $s=1+\eps$.
The reciprocal-Gamma expansion gives
\[
 \frac{\tau^{1+\eps}}{\Gamma(1+\eps)}
   =\tau\sum_{j\ge0}\Ecal_j(L)\frac{\eps^j}{j!},
 \qquad \Bcal_0(1+\eps;a,\tau)=\Qcal_{-\eps}(a).
\]
The normally convergent lattice series permits coefficient extraction
using \eqref{eq:K-generator}. The Leibniz coefficient is
$(-1)^m(-1)^{m-j}=(-1)^j$, giving the stated sign. All order
extractions on the right are permitted by \eqref{eq:Lambert}.
\end{proof}

\subsection{The first three identities}
With $\Qcal_0^{[r]}=\left.\partial_\nu^r\Qcal_\nu\right|_{\nu=0}$,
\begin{align}
 \sum_n\left(\Log w_n-\psi(w_n)-\frac1{2w_n}\right)
   &=\tau\Qcal_0,\label{eq:K0}\\
 \sum_n\left(\gamma_1(w_n)+\frac{\Log^2w_n}{2}
                         -\frac{\Log w_n}{2w_n}\right)
   &=\tau(\Qcal_0^{[1]}-L\Qcal_0),\label{eq:K1}\\
 \sum_n\left(\gamma_2(w_n)+\frac{\Log^3w_n}{3}
                         -\frac{\Log^2w_n}{2w_n}\right)
   &=\tau\bigl(\Qcal_0^{[2]}-2L\Qcal_0^{[1]}
                     +(L^2-\zeta(2))\Qcal_0\bigr).
 \label{eq:K2}
\end{align}
For real $a$, opposite lattice terms are conjugate; the left sides
are real. The right sides are real as well, since $0<z,q<1$.
For example, at $a=0.8$ and $h=2.1$, \eqref{eq:K0} is approximately
$0.0690653783784057$. This displayed decimal is a diagnostic value,
not an interval enclosure.

\subsection{Finite horizontal differences remove the Lambert series}
For an integer $J\ge1$, geometric cancellation gives
\begin{equation}\label{eq:Q-finite}
 \begin{split}
 \Qcal_\nu(a)-\Qcal_\nu(a+J)
 ={}&\sum_{r=0}^{J-1}\Li_\nu(zq^r)
       -\frac12\bigl(\Li_\nu(z)-\Li_\nu(zq^J)\bigr)\\
    &-\frac1\tau\bigl(\Li_{\nu+1}(z)-\Li_{\nu+1}(zq^J)\bigr).
 \end{split}
\end{equation}
Thus the difference of the two Stieltjes lattice sums in
\eqref{eq:stieltjes-main} has a \emph{finite} polylogarithmic
expression, with all order derivatives and Bell coefficients explicit.
In particular, for $J=1$ the bracket is
\begin{equation}\label{eq:Q-unit}
 \frac12\bigl(\Li_\nu(z)+\Li_\nu(zq)\bigr)
 -\frac1\tau\bigl(\Li_{\nu+1}(z)-\Li_{\nu+1}(zq)\bigr).
\end{equation}
Equations \eqref{eq:stieltjes-main} and \eqref{eq:Q-unit} constitute
an infinite family of bilateral generalized-Stieltjes identities with
only finitely many polylogarithms on the right. Their derivation uses
only absolute sums. They do not make an arithmetic independence claim
about the order derivatives that occur.

\section{Differentiation and normalized antiderivatives}
\label{sec:calculus}

The Mellin normalization makes horizontal calculus transparent:
\begin{equation}\label{eq:horizontal-derivative}
 \partial_w^r\Rcal_M(s,w)=(-1)^r(s)_r\Rcal_M(s+r,w).
\end{equation}
This follows by differentiation of \eqref{eq:mellin} in a common
convergence domain and then analytic continuation. Accordingly,
whenever the original lattice sum lies in its normal domain,
\begin{equation}\label{eq:lattice-derivative}
 \partial_a^r\Scal_M(s;a,h)=(-1)^r(s)_r\Scal_M(s+r;a,h).
\end{equation}
Subsequent spectral differentiation is legitimate on every compact
subset of that domain. The derivative of the rising factorial is
part of the formula and must not be omitted.

There is an equally simple integral form. Since
\begin{equation}\label{eq:Q-derivative}
 \partial_a\Qcal_\nu(a)=-\tau\Qcal_{\nu-1}(a),
\end{equation}
raising the order supplies an anchored primitive.

\begin{theorem}[All-index anchored primitives]\label{thm:primitives}
Let $a_0,a$ lie in the right half-plane and integrate along any path
there. Then
\begin{equation}\label{eq:primitive-stieltjes}
 \begin{split}
 \int_{a_0}^a\sum_{n\in\Z}K_m(u+\ii hn)\dd u
 ={}&-\sum_{j=0}^m\binom mj(-1)^j\Ecal_j(L)\\
 &\quad\cdot\left.
 \partial_\nu^{m-j}\bigl(\Qcal_\nu(a)-\Qcal_\nu(a_0)\bigr)
 \right|_{\nu=1}.
 \end{split}
\end{equation}
In particular, the primitive vanishes at $a=a_0$. More generally,
\begin{equation}\label{eq:primitive-spectral}
 \int_{a_0}^a\Scal_M(s;u,h)\dd u
 =-\frac{\tau^{s-1}}{\Gamma(s)}
       \bigl(\Bcal_M(s-1;a,\tau)-\Bcal_M(s-1;a_0,\tau)\bigr)
\end{equation}
on the normal domain $\Ree s>-2M$.
\end{theorem}
\begin{proof}
Differentiate the right sides using
$\partial_a\Li_\nu(e^{-\tau a})=-\tau\Li_{\nu-1}(e^{-\tau a})$
and its normally convergent Lambert counterpart. This gives the
integrands by Theorems~\ref{thm:sampling} and \ref{thm:stieltjes}.
The expressions vanish at $a_0$, and the right half-plane is simply
connected. Normal convergence justifies the integral of the lattice
sum along a compact path. This proof avoids an unjustified attempt
to integrate each nonconvergent primitive separately.
\end{proof}

When $\Ree(s-1)>-2M$ and $s\ne1$, \eqref{eq:primitive-spectral}
may equivalently be written
\[
 \frac{\Scal_M(s-1;a_0,h)-\Scal_M(s-1;a,h)}{s-1}.
\]
Outside that smaller domain, this last notation refers to the
continued summed function, not necessarily an ordinary series of
primitives. Formula~\eqref{eq:primitive-spectral} itself remains
unambiguous on the larger stated domain.

\section{The Gamma specialization and its known product identity}
\label{sec:gamma}

Define the normalized Stirling remainder
\begin{equation}\label{eq:G}
 G(w)=\log\Gamma(w)-(w-\tfrac12)\Log w+w-\tfrac12\log(2\pi).
\end{equation}
The Hurwitz special value $\zeta'(0,w)=\log\Gamma(w)-\tfrac12\log(2\pi)$
implies
\begin{equation}\label{eq:Gamma-jet}
 \partial_s\Rcal_1(0,w)=G(w)-\frac1{12w}.
\end{equation}
This is inside the normal domain of $\Scal_1$. Differentiating
\eqref{eq:master} there gives the absolutely convergent identity
\begin{equation}\label{eq:Gamma-absolute}
 \sum_{n\in\Z}\left(G(w_n)-\frac1{12w_n}\right)
 =\Lcal_1(q,z)-\frac12\Li_1(z)-\frac1\tau\Li_2(z)
                         -\frac\tau{12}\Li_0(z).
\end{equation}

For a function with reciprocal decay we specify
\[
 \Sym\sum_{n\in\Z}F(w_n):=\lim_{N\to\infty}\sum_{n=-N}^{N}F(w_n).
\]
The elementary partial-fraction identity is
\begin{equation}\label{eq:reciprocal-sym}
 \Sym\sum_{n\in\Z}\frac1{a+\ii hn}
     =\frac\tau2\coth\frac{\tau a}{2}
     =\tau\left(\Li_0(z)+\frac12\right).
\end{equation}
For completeness, pair $n$ and $-n$ to obtain
$1/a+\sum_{n\ge1}2a/(a^2+h^2n^2)$; the standard partial fractions
of $\coth$ give \eqref{eq:reciprocal-sym}. They follow as well by
logarithmically differentiating the canonical product for $\sinh$.

\begin{corollary}[Known Gamma-product specialization]
\label{cor:Gamma}
For $h>0$ and $\Ree a>0$,
\begin{equation}\label{eq:Gamma-symmetric}
 \boxed{\displaystyle
 \Sym\sum_{n\in\Z}G(a+\ii hn)
  =-\log(z;q)_\infty+\frac12\log(1-z)
       -\frac1\tau\Li_2(z)+\frac\tau{24}.}
\end{equation}
Here $\log(z;q)_\infty=\sum_{r\ge0}\log(1-zq^r)$, with the
analytic logarithm supplied by its convergent series.
\end{corollary}
\begin{proof}
Add one twelfth of \eqref{eq:reciprocal-sym} to
\eqref{eq:Gamma-absolute}, use
$\Lcal_1=-\log(z;q)_\infty$ and $\Li_1(z)=-\log(1-z)$.
The $\Li_0$ terms cancel, leaving $\tau/24$.
\end{proof}

Equation~\eqref{eq:Gamma-symmetric}, after exponentiation and the
substitution $y=\tau a$, $\beta=\tau$, is precisely the Gamma-product
specialization credited to Arabi Ardehali and Rosengren
\cite[Eq.~(13)]{ArdehaliRosengren}. Its absolute version corresponds
to their subtracted remainder. We rederive it only to fix the
normalization of the subsequent general contact law; neither the
Gamma product nor its constant is claimed as new.

There is a useful interchange diagnostic. Differentiating the entire
continuation of $\Scal_0(s)$ at $s=0$ gives
$\Bcal_0(0)$, whereas symmetrically summing
$\partial_s\Rcal_0(0,w_n)=G(w_n)$ gives
$\Bcal_0(0)+\tau/24$. The discrepancy is exactly $\tau/24$, not a
numerical error. At $a=0.8$, $h=2.1$, the correct symmetric value is
approximately $0.14613043538142660274$.

\section{Negative integer jets: a parity-sensitive contact law}
\label{sec:negative}

For $N\ge0$, put
\begin{equation}\label{eq:T-definition}
 M_N=\left\lceil\frac N2\right\rceil,
 \qquad T_N(w)=\left.\partial_s\Rcal_{M_N}(s,w)\right|_{s=-N}.
\end{equation}
The chosen compensation is minimal for the parity statement below.
Every $\Rcal_{M_N}(-N,w)$ is zero. This follows from
\eqref{eq:mellin}, whose integral is convergent there, or directly
from the Bernoulli-polynomial formula for $\zeta(-N,w)$.

\begin{proposition}[Explicit harmonic-number form]\label{prop:T-explicit}
For $N\ge0$ and $\Ree w>0$,
\begin{equation}\label{eq:T-explicit}
 \begin{split}
 T_N(w)={}&\zeta'(-N,w)-\frac{w^{N+1}\Log w}{N+1}
                   +\frac{w^{N+1}}{(N+1)^2}
                   +\frac12w^N\Log w\\
 &-\sum_{j=1}^{M_N}\frac{B_{2j}}{(2j)!}
    \frac{N!}{(N-2j+1)!}\,w^{N-2j+1}
    \bigl(\Log w+H_N-H_{N-2j+1}\bigr).
 \end{split}
\end{equation}
The sum is empty when $N=0$.
\end{proposition}
\begin{proof}
For $k=2j-1\le N$,
\[
 (-N)_k=-\frac{N!}{(N-k)!},\qquad
 \left.\frac{\partial_s(s)_k}{(s)_k}\right|_{s=-N}
       =-H_N+H_{N-k}.
\]
Differentiate each term of \eqref{eq:remainder}. The two minus signs
in the derivative of $(s)_k w^{1-s-2j}$ give the displayed
harmonic factor. Differentiating the polar subtraction contributes
$-w^{N+1}\Log w/(N+1)+w^{N+1}/(N+1)^2$.
\end{proof}

\begin{lemma}[The first nonzero reciprocal coefficient]\label{lem:beta}
If $N$ is even, then
\begin{equation}\label{eq:T-even-asymptotic}
 T_N(w)=\frac{\beta_N}{w}+O(w^{-3}),\qquad
 \beta_N=\frac{B_{N+2}}{(N+1)(N+2)}.
\end{equation}
If $N$ is odd, then $T_N(w)=O(w^{-2})$.
These estimates are uniform on closed sectors avoiding the negative axis.
\end{lemma}
\begin{proof}
For $k\ge N+1$, the product $(s)_k$ has exactly one zero factor at
$s=-N$, and
\begin{equation}\label{eq:rf-zero-derivative}
 \left.\partial_s(s)_k\right|_{s=-N}
       =(-1)^N N!(k-N-1)!.
\end{equation}
Differentiate the asymptotic expansion
\eqref{eq:em-asymptotic}. No logarithm remains in its terms because
$(s)_k$ itself vanishes at $-N$. The first remaining exponent is
$-1$ for even $N$ and $-2$ for odd $N$. In the even case its
coefficient is \eqref{eq:T-even-asymptotic}; successive exponents
differ by two.
\end{proof}

\begin{theorem}[All nonpositive integer first jets]\label{thm:negative}
For $N\ge0$, $h>0$, and $\Ree a>0$,
\begin{equation}\label{eq:negative-main}
 \boxed{\displaystyle
 \Sym\sum_{n\in\Z}T_N(a+\ii hn)
 =(-1)^N N!\tau^{-N}\Bcal_{M_N}(-N;a,\tau)
 +\mathbf1_{2\mid N}\frac{\tau B_{N+2}}{2(N+1)(N+2)}.}
\end{equation}
For odd $N$, the sum is absolutely convergent and $\Sym$ may be
omitted. For even $N$, the displayed symmetric prescription exists;
the contact is independent of $a$.
\end{theorem}
\begin{proof}
When $N$ is odd, $-N>-2M_N$. Normal convergence in
Theorem~\ref{thm:sampling} permits differentiation at $s=-N$.
Since $(1/\Gamma)'(-N)=(-1)^N N!$, this gives the first term on
the right and no contact.

Let $N=2M$ be even. Oversubtract once, using $\Rcal_{M+1}$.
The order $-N$ is now strictly inside the normal convergence domain.
The difference of the two first jets is exactly $\beta_N/w$, by
\eqref{eq:rf-zero-derivative}. Thus
\[
 \Sym\sum_n T_N(w_n)
 =N!\tau^{-N}\Bcal_{M+1}(-N;a,\tau)
        +\beta_N\frac\tau2\coth\frac{\tau a}{2}.
\]
By the definition of $\Bcal_M$,
\[
 N!\tau^{-N}\Bcal_{M+1}(-N)
 =N!\tau^{-N}\Bcal_M(-N)-\beta_N\tau\Li_0(z).
\]
Use $\coth(\tau a/2)=1+2\Li_0(z)$. The image terms cancel and
leave $\beta_N\tau/2$. Existence of the symmetric sum follows
from Lemma~\ref{lem:beta} and \eqref{eq:reciprocal-sym}.
\end{proof}

\subsection{Two explicit higher identities}
The first negative odd order has
\begin{align}
 T_1(w)={}&\zeta'(-1,w)-\tfrac12w^2\Log w+\tfrac14w^2
                 +\tfrac12w\Log w-\tfrac1{12}(\Log w+1),\label{eq:T1}\\
 \sum_{n\in\Z}T_1(w_n)
   ={}&-\frac{\Lcal_2(q,z)}\tau+\frac{\Li_2(z)}{2\tau}
                    +\frac{\Li_3(z)}{\tau^2}+\frac{\Li_1(z)}{12}.
 \label{eq:T1-sum}
\end{align}
This is an ordinary absolutely convergent identity. The next order is
\begin{align}
 T_2(w)={}&\zeta'(-2,w)-\tfrac13w^3\Log w+\tfrac19w^3
               +\tfrac12w^2\Log w-\tfrac16w\Log w-\tfrac1{12}w,
 \label{eq:T2}\\
 \Sym\sum_{n\in\Z}T_2(w_n)
   ={}&\frac{2}{\tau^2}\Lcal_3(q,z)-\frac{\Li_3(z)}{\tau^2}
             -\frac{2\Li_4(z)}{\tau^3}-\frac{\Li_2(z)}{6\tau}
             -\frac\tau{720}.
 \label{eq:T2-sum}
\end{align}
The final $-\tau/720$ is required. These examples display explicitly
how negative-order Hurwitz derivatives connect to polylogarithms of
successive positive integer orders. Formula~\eqref{eq:T-explicit}
provides the corresponding harmonic-number expression at every order.

\subsection{Periodic filters annihilate the mean contact}
\begin{corollary}[Mean-zero periodic weights]\label{cor:filters}
Let $c_n$ be periodic of period $d\ge1$. Put
$\tau_d=\tau/d$ and $a_r=a+\ii hr$ for $0\le r<d$. Then
\begin{equation}\label{eq:filtered}
 \begin{split}
 \Sym\sum_{n\in\Z}c_n T_N(a+\ii hn)
 ={}&(-1)^N N!\tau_d^{-N}
       \sum_{r=0}^{d-1}c_r\Bcal_{M_N}(-N;a_r,\tau_d)\\
 &+\mathbf1_{2\mid N}\frac{\tau_d B_{N+2}}{2(N+1)(N+2)}
                  \sum_{r=0}^{d-1}c_r.
 \end{split}
\end{equation}
Consequently every first-jet contact vanishes for a mean-zero filter.
\end{corollary}
\begin{proof}
Split the lattice into residue classes modulo $d$ and apply
Theorem~\ref{thm:negative} to each class with spacing $dh$.
For even $N$, shifting the endpoints of a symmetric cutoff by a
bounded number of terms changes its limit by zero, because
$T_N(w_n)=O(1/|n|)$. For odd $N$ regrouping is absolute.
The contact in each class is the same constant. Summing yields
\eqref{eq:filtered}.
\end{proof}

For example, the alternating Gamma specialization is
\begin{equation}\label{eq:alternating-Gamma}
 \Sym\sum_{n\in\Z}(-1)^nG(a+\ii hn)
      =2\Bcal_0(0;a,\tau/2)-\Bcal_0(0;a,\tau).
\end{equation}
Indeed the even-index lattice has spacing $2h$, and the two contacts
cancel. Formula~\eqref{eq:filtered} identifies the obstruction as a
single mean mode, not an arbitrary choice of regularization.

\section{A harmonic-number compiler for every higher spectral jet}
\label{sec:all-jets}

At a negative integer, reciprocal Gamma has the exact local factorization
\begin{equation}\label{eq:rgamma-factor}
 \frac{\tau^{-N+\eps}}{\Gamma(-N+\eps)}
 =(-1)^N N!\tau^{-N}\eps\,
   \frac{\tau^\eps}{\Gamma(1+\eps)}
   \prod_{k=1}^N\left(1-\frac\eps k\right).
\end{equation}
Define $\Ecal_j^{(N)}$ by
\begin{equation}\label{eq:EN}
 \begin{split}
 &\exp\left((\log\tau+\gamma-H_N)u
   +\sum_{r=2}^\infty
     \frac{(-1)^{r+1}\zeta(r)-H_N^{(r)}}r u^r\right)\\
 &\hspace{30mm}=\sum_{j=0}^\infty\Ecal_j^{(N)}\frac{u^j}{j!}.
 \end{split}
\end{equation}
For $L_N=\log\tau+\gamma-H_N$,
\begin{align*}
 \Ecal_0^{(N)}&=1,\qquad \Ecal_1^{(N)}=L_N,\\
 \Ecal_2^{(N)}&=L_N^2-\zeta(2)-H_N^{(2)},\\
 \Ecal_3^{(N)}&=L_N^3-3L_N(\zeta(2)+H_N^{(2)})
                         +2\zeta(3)-2H_N^{(3)}.
\end{align*}
These formulas keep the generalized harmonic correction separate from
the alternating sign of the zeta contribution.

\begin{theorem}[Every compensated negative-integer jet]
\label{thm:negative-jets}
Let $N\ge0$, $L$ be an integer with $L>N/2$, and $d\ge1$.
Then the following lattice sum is absolutely convergent:
\begin{equation}\label{eq:negative-jets}
 \begin{split}
 &\sum_{n\in\Z}
       \left.\partial_s^d\Rcal_L(s,a+\ii hn)\right|_{s=-N}\\
 &\quad=(-1)^N N!\tau^{-N}d
       \sum_{j=0}^{d-1}\binom{d-1}{j}\Ecal_j^{(N)}
           \left.\partial_s^{d-1-j}\Bcal_L(s;a,\tau)\right|_{s=-N}.
 \end{split}
\end{equation}
Every derivative on the right is an explicit polylogarithm order
derivative or a Lambert sum of such derivatives.
\end{theorem}
\begin{proof}
The strict inequality $L>N/2$ places $-N$ inside the normal
convergence domain of Theorem~\ref{thm:sampling}. Expand its right
side using \eqref{eq:rgamma-factor}. Taking the logarithm of the
finite product gives $-\sum_{r\ge1}H_N^{(r)}\eps^r/r$, which,
together with the reciprocal-Gamma series, yields \eqref{eq:EN}.
The coefficient of $\eps^d$ in $\eps\Ecal^{(N)}(\eps)
\Bcal_L(-N+\eps)$ is a Leibniz coefficient of degree $d-1$.
Multiplication by $d!$ gives exactly the factor
$d\binom{d-1}{j}$ in \eqref{eq:negative-jets}.
\end{proof}

This theorem is a finite symbolic procedure at every prescribed
$(N,d,L)$. It involves no numerical guessing of coefficients.
For the minimally compensated even-order first jet one may instead
use Theorem~\ref{thm:negative}; for higher minimal jets, logarithmic
reciprocal tails generally appear. No ordinary-sum interpretation of
those insufficiently compensated higher jets is asserted here.

\section{Products of polygammas: finite Hurwitz--Lerch images}
\label{sec:bilinear}

In this section $a,b>0$ are real. Let $r,t\ge2$ be integers, set
$A=a+b$, and define
\begin{equation}\label{eq:D}
 D_\ell(Q,A)=\frac{\zeta(\ell,A)-Q\Phi(Q,\ell,A)}{1-Q},
 \qquad 0<Q<1,\quad \ell\ge2,
\end{equation}
where
\[
 \Phi(Q,\ell,A)=\sum_{m=0}^\infty\frac{Q^m}{(m+A)^\ell}.
\]
The elementary partial fractions
\begin{equation}\label{eq:denominator-pf}
 \frac1{(1-x)(1-Qx)}
   =\frac1{1-Q}\left(\frac1{1-x}-\frac{Q}{1-Qx}\right)
\end{equation}
give the convergent integral
\begin{equation}\label{eq:D-integral}
 \int_0^\infty\frac{u^{\ell-1}e^{-Au}}
               {(1-e^{-u})(1-Qe^{-u})}\dd u
        =\Gamma(\ell)D_\ell(Q,A).
\end{equation}

\begin{theorem}[Finite-image product identity]\label{thm:bilinear}
For $a,b,h>0$ and integers $r,t\ge2$,
\begin{equation}\label{eq:bilinear-main}
 \sum_{n\in\Z}\zeta(r,a+\ii hn)\zeta(t,b-\ii hn)
       =V_{r,t}(A)+\sum_{k=1}^\infty
             \bigl(I_{r,t;k}(a,b)+I_{t,r;k}(b,a)\bigr),
\end{equation}
where
\begin{align}
 V_{r,t}(A)
  &=\frac{\tau\Gamma(r+t-1)}{\Gamma(r)\Gamma(t)}
       \bigl(\zeta(r+t-2,A)-(A-1)\zeta(r+t-1,A)\bigr),
       \label{eq:V}\\
 I_{r,t;k}(a,b)
  &=\frac{\tau q^{ak}}{\Gamma(r)\Gamma(t)}
       \sum_{j=0}^{r-1}\binom{r-1}{j}(\tau k)^{r-1-j}
                      \Gamma(j+t)D_{j+t}(q^k,A).
       \label{eq:I}
\end{align}
The lattice sum is absolutely convergent. The image series converges
exponentially, and each individual image is a finite combination of
Hurwitz zeta and Lerch values.
\end{theorem}
\begin{proof}
Extend the two functions
\[
 f(u)=\frac{u^{r-1}e^{-au}}{\Gamma(r)(1-e^{-u})},\qquad
 g(u)=\frac{u^{t-1}e^{-bu}}{\Gamma(t)(1-e^{-u})}
\]
by zero on the negative axis. They are in $L^1\cap L^2$.
Their Fourier transforms are $\zeta(r,a+\ii y)$ and
$\zeta(t,b+\ii y)$. The correlation
\[
 P(x)=\int_{\mathbb R}f(v+x)g(v)\dd v
\]
is continuous by continuity of translations in $L^2$, decays
exponentially up to polynomial factors in both directions, and has
Fourier transform
$\widehat f(y)\widehat g(-y)$. The sampled transform is absolutely
summable since the product is $O(|n|^{2-r-t})$.
Apply Lemma~\ref{lem:periodization} to $P$.

At $x=0$ expand $(1-e^{-u})^{-2}=\sum_{m\ge0}(m+1)e^{-mu}$.
Integration gives \eqref{eq:V}, since
$m+1=(m+A)-(A-1)$. For $x=\tau k>0$ expand
$(v+\tau k)^{r-1}$ by the binomial theorem and apply
\eqref{eq:D-integral}. This gives \eqref{eq:I}. Negative $x$
swaps $(r,a)$ and $(t,b)$. Exponential decay of the images follows
also directly from the integral, with factors $q^{ak}$ or $q^{bk}$
and a fixed polynomial in $k$.
\end{proof}

The usual relation
\begin{equation}\label{eq:polygamma}
 \zeta(r,w)=\frac{(-1)^r}{(r-1)!}\psi^{(r-1)}(w)
\end{equation}
converts the theorem into a formula for every product of two positive
polygamma derivatives. It also permits a finite harmonic-number
interpretation after horizontal integer differences are taken.

\subsection{An explicit trigamma square}
For $a>0$, $r=t=2$ and $b=a$ give
\begin{equation}\label{eq:trigamma-square}
 \begin{split}
 \sum_{n\in\Z}|\psi'(a+\ii hn)|^2
 ={}&2\tau\bigl(\zeta(2,2a)-(2a-1)\zeta(3,2a)\bigr)\\
 &+2\tau\sum_{k=1}^\infty q^{ak}
       \bigl(\tau kD_2(q^k,2a)+2D_3(q^k,2a)\bigr).
 \end{split}
\end{equation}
The $k=0$ term is the continuous vertical pairing at the sampling
normalization. Every $k\ge1$ term is a discrete image correction.
Keeping only the continuous term is therefore not an exact lattice
identity.

\subsection{Finite generalized-harmonic differences}
The shift formula
\[
 \zeta(r,w)-\zeta(r,w+J)=\sum_{j=0}^{J-1}(w+j)^{-r}
\]
turns four horizontal evaluations of \eqref{eq:bilinear-main} into
an exact pairing of finite generalized-harmonic sums. All finite sums
may be interchanged with the lattice series because its terms decay
at least quadratically. Thus the same identity system covers both
infinite polygamma tails and finite harmonic windows; no separate
regularization is needed.

\section{All-index Stieltjes differences and four-shift closure}
\label{sec:differences}

Let $a,c,b,d>0$ and write
\begin{align}
 \Delta_{a,c}(s;n)&=\zeta(s,a+\ii hn)-\zeta(s,c+\ii hn),
 \label{eq:Delta}\\
 \delta_m^{a,c}(n)&=\gamma_m(a+\ii hn)-\gamma_m(c+\ii hn).
 \label{eq:delta}
\end{align}
The pole at $s=1$ cancels in $\Delta$, and
\[
 \Delta_{a,c}(1+\eps;n)
    =\sum_{m\ge0}\frac{(-1)^m}{m!}\delta_m^{a,c}(n)\eps^m.
\]
For each fixed $m$, $\delta_m^{a,c}(n)=O(\log^m(2+|n|)/(1+|n|))$.
Its product with a second such difference is absolutely summable.

For $\Ree s>0$ set
\begin{equation}\label{eq:diff-density}
 f_s^{a,c}(u)=\frac{u^{s-1}}{\Gamma(s)}
                 \frac{e^{-au}-e^{-cu}}{1-e^{-u}},\qquad u>0,
\end{equation}
extended by zero to $u<0$. It belongs to $L^1$; when $\Ree s>1/2$
it also belongs to $L^2$. Its Fourier transform is
$\Delta_{a,c}(s;y/h)$ in the evident continuous-variable notation.

\begin{theorem}[Difference pairing and all its Stieltjes coefficients]
\label{thm:diff-jets}
For $\Ree s,\Ree t>1/2$,
\begin{equation}\label{eq:diff-pair}
 \begin{split}
 &\sum_{n\in\Z}\Delta_{a,c}(s;n)\Delta_{b,d}(t;-n)\\
 &=\tau\int_0^\infty f_s^{a,c}(u)f_t^{b,d}(u)\dd u\\
 &\quad+\tau\sum_{k=1}^\infty\int_0^\infty
   \bigl(f_s^{a,c}(u+\tau k)f_t^{b,d}(u)
        +f_s^{a,c}(u)f_t^{b,d}(u+\tau k)\bigr)\dd u.
 \end{split}
\end{equation}
All series, integrals, and fixed spectral derivatives converge
normally near $(s,t)=(1,1)$. Define
\begin{equation}\label{eq:Stieltjes-density}
 F_m^{a,c}(u)=(-1)^m\Ecal_m(\log u+\gamma)
                        \frac{e^{-au}-e^{-cu}}{1-e^{-u}}.
\end{equation}
Then, for every $m,\ell\ge0$,
\begin{equation}\label{eq:all-diff-jets}
 \begin{split}
 \sum_{n\in\Z}\delta_m^{a,c}(n)\delta_\ell^{b,d}(-n)
 ={}&\tau\int_0^\infty F_m^{a,c}(u)F_\ell^{b,d}(u)\dd u\\
 &+\tau\sum_{k\ge1}\int_0^\infty
   \bigl(F_m^{a,c}(u+\tau k)F_\ell^{b,d}(u)\\
 &\hspace{36mm}+F_m^{a,c}(u)F_\ell^{b,d}(u+\tau k)\bigr)\dd u.
 \end{split}
\end{equation}
These are ordinary absolutely convergent identities.
\end{theorem}
\begin{proof}
Use the correlation and periodization argument of
Theorem~\ref{thm:bilinear}, now for the $L^1\cap L^2$ densities
\eqref{eq:diff-density}. The shift difference improves the vertical
Hurwitz decay to $O(|n|^{-\Ree s})$, so the sampled product is
absolutely summable in the stated domain. At zero, the densities
behave as a constant times $u^{s-1}$; at infinity they decay
exponentially. These facts give the required normal convergence,
including logarithmic factors from spectral derivatives.

The coefficient identity
$(-1)^m\partial_s^m f_s^{a,c}|_{s=1}=F_m^{a,c}$ follows from
\eqref{eq:E}. Extract the two coefficients in \eqref{eq:diff-pair}.
The signs and factorials then match \eqref{eq:all-diff-jets}.
\end{proof}

\subsection{The central image has finite Hurwitz--Stieltjes coordinates}
Let $a_1=a$, $a_2=c$, $b_1=b$, $b_2=d$, and
$\epsilon_1=1$, $\epsilon_2=-1$. Put
\begin{equation}\label{eq:Hcentral}
 \mathcal H(v)=\sum_{i,j=1}^2\epsilon_i\epsilon_j
       \bigl(\zeta(v-1,a_i+b_j)
                  -(a_i+b_j-1)\zeta(v,a_i+b_j)\bigr).
\end{equation}
The potential poles at $v=2$ and $v=1$ cancel because
\[
 \sum_{i,j}\epsilon_i\epsilon_j=0,
 \qquad \sum_{i,j}\epsilon_i\epsilon_j(a_i+b_j-1)=0.
\]
Hence $\mathcal H$ is entire. The central term of
\eqref{eq:diff-pair} is
\begin{equation}\label{eq:central-spectral}
 \tau\frac{\Gamma(s+t-1)}{\Gamma(s)\Gamma(t)}\mathcal H(s+t-1).
\end{equation}
To prove it, expand the squared geometric denominator in a common
absolute-convergence domain and integrate termwise, then continue
the grouped expression to $\Ree s,\Ree t>1/2$.

In particular,
\begin{equation}\label{eq:central-jets}
 \mathcal H^{(k)}(1)=\sum_{i,j}\epsilon_i\epsilon_j
 \left(\zeta^{(k)}(0,a_i+b_j)
       -(-1)^k(a_i+b_j-1)\gamma_k(a_i+b_j)\right).
\end{equation}
Every central Stieltjes coefficient is therefore a finite
combination of these values and reciprocal-Gamma Bell coefficients.
Explicitly, the central contribution to \eqref{eq:all-diff-jets} is
\begin{equation}\label{eq:central-coeff}
 \tau(-1)^{m+\ell}m!\ell!
 [u^m v^\ell]
 \frac{\Gamma(1+u+v)}{\Gamma(1+u)\Gamma(1+v)}
                          \mathcal H(1+u+v).
\end{equation}
The nonzero images in \eqref{eq:all-diff-jets} remain explicit
one-dimensional logarithmic integrals. We do not claim a finite
single-Hurwitz closure for arbitrary spectral derivatives of those
images.

\subsection{A fully explicit four-shift digamma identity}
\begin{theorem}[Four-shift closure]\label{thm:fourshift}
With the preceding positive real parameters and
$A_{ij}=a_i+b_j$, define
\[
 \mathcal P=\sum_{n\in\Z}
   \bigl(\psi(c+\ii hn)-\psi(a+\ii hn)\bigr)
   \bigl(\psi(d-\ii hn)-\psi(b-\ii hn)\bigr).
\]
Then
\begin{equation}\label{eq:four-shift}
 \begin{split}
 \mathcal P={}&\tau\sum_{i,j}\epsilon_i\epsilon_j
                             (A_{ij}-1)\psi(A_{ij})\\
 &-\tau\sum_{k=1}^\infty\frac1{1-q^k}
     \sum_{i,j}\epsilon_i\epsilon_j
       \bigl(q^{a_i k}+q^{b_j k}\bigr)\\
 &\hspace{34mm}\cdot
        \bigl(\psi(A_{ij})+q^k\Phi(q^k,1,A_{ij})\bigr).
 \end{split}
\end{equation}
The original lattice series is absolutely convergent and the image
series converges exponentially.
\end{theorem}
\begin{proof}
Put $s=t=1$ in \eqref{eq:diff-pair}. The central image is
$\tau\mathcal H(1)$; the $\zeta(0,A)$ terms in
\eqref{eq:central-jets} cancel, giving the first line of
\eqref{eq:four-shift}. For a positive image, use
\eqref{eq:denominator-pf}. Its would-be Hurwitz order-one pole is
independent of $A_{ij}$, and its total coefficient is
\[
 \sum_{i,j}\epsilon_i\epsilon_j q^{a_i k}=0.
\]
Thus the finite part of the grouped expression uses
$\zeta(1+\eta,A)=\eta^{-1}-\psi(A)+O(\eta)$ and gives
\[
 -\frac{\tau}{1-q^k}\sum_{i,j}\epsilon_i\epsilon_j q^{a_i k}
            \bigl(\psi(A_{ij})+q^k\Phi(q^k,1,A_{ij})\bigr).
\]
This is an ordinary convergent grouped integral; the Laurent
calculation only evaluates it. The negative image has $q^{b_jk}$
instead. Adding them proves \eqref{eq:four-shift}.
\end{proof}

\begin{example}[A normalization check that reduces to elementary functions]
If $c=a+1$ and $d=b+1$, the two digamma differences are reciprocals.
Consequently \eqref{eq:four-shift} must reduce to
\begin{equation}\label{eq:fourshift-control}
 \sum_{n\in\Z}\frac1{(a+\ii hn)(b-\ii hn)}
 =\frac\tau{2(a+b)}
       \left(\coth\frac{\tau a}{2}+\coth\frac{\tau b}{2}\right).
\end{equation}
This follows independently by partial fractions and
\eqref{eq:reciprocal-sym}. The paired product is itself absolutely
convergent. Longer integer shift windows reduce to finite sums of
\eqref{eq:fourshift-control}, giving further independent checks of
all signs and scale factors in the four-shift image formula.
\end{example}

\section{Verification, corrections, and integration status}
\label{sec:audit}

\subsection{What was checked}
The distributed \src{verification/run_checks.py} executes 268 exact
algebraic assertions and 38 floating-point comparisons at 65 working
decimal digits. All passed in the recorded run. The largest recorded
scaled residual was approximately $3.55\times10^{-43}$; tolerances
were $10^{-32}$ or $10^{-30}$, depending on the test. These residuals
are diagnostics, not rigorous error bounds.

The exact assertions verify Bernoulli-polynomial cancellation at
negative integer orders, the zero-factor derivative
\eqref{eq:rf-zero-derivative}, the first contact coefficient,
harmonic-number factors in \eqref{eq:T-explicit}, and the rational
partial fractions used by the image formulas. They test finitely
many instances of the uniform proofs; they are not presented as a
substitute for those proofs.

The numerical routes are deliberately separated. The lattice side
uses direct Hurwitz evaluations for a finite central range, then a
sectorial Euler--Maclaurin tail summed through Hurwitz zeta values.
The Lambert side uses rapidly convergent real-exponential series.
Stieltjes coefficients are obtained from a pole-cancelled,
complex-safe Euler--Maclaurin expression rather than a call to
\src{mpmath.stieltjes} at nonreal arguments. Integer product
pairings use independently convolved large-lattice asymptotic
coefficients. Four-shift tests include finite reciprocal reductions
and direct integral checks of noninteger-shift images.

Some checks compare two symbolic consequences of the same master
identity, notably the higher-jet Bell compiler and primitive
transport; these are labeled as cross-route checks, not as fully
independent validations of the master theorem. Two corruption
controls reject a changed zero-factor derivative and omission of
the Gamma contact. Complete numerical values, tolerances, versions,
and classifications are recorded in \src{results/verification.json}.

\subsection{A concrete editorial correction}
The inspected manuscript README contains the same paragraph beginning
``The complete 1,032-input first-letter X family'' twice in immediate
succession \cite{RepoMain}. The proposed correction is to retain one
copy and delete the duplicate, without changing any mathematical
status statement. This is an editorial correction, not evidence of
a false theorem. The package contains a guarded script that removes
the second occurrence only when the exact duplicated paragraph is
present in a supplied local README copy. It does not contact or
modify GitHub.

\subsection{A mathematical warning, not an attributed repository error}
The tempting identity
\[
 \Sym\sum_n G(w_n)\stackrel{\mathrm{wrong}}=
     \left.\frac{\partial}{\partial s}
          \frac{\tau^s\Bcal_0(s)}{\Gamma(s)}\right|_{s=0}
\]
omits $\tau/24$. More generally, the omitted term at even negative
order $-N$ is the contact in \eqref{eq:negative-main}. We provide
this as an explicit correction to that invalid interchange. The
audit did \emph{not} establish that an inspected canonical theorem
made this mistake, and it should not be recorded as an error in a
particular existing manuscript without an actual source occurrence.

\subsection{Placement and proof status}
A natural placement is a new continuation under the existing
vertical-pairing report. The core depends only on the standard
Hurwitz/Mellin identities and the elementary periodization lemma.
It does not require accepting any pending continuum or arithmetic
claim in an incoming archive. A proposed integration order is:
compensation and sampling; Stieltjes jets and primitives;
negative-integer contacts; higher jets; integer pairings;
Stieltjes differences and four-shift closure.

All theorem statements in this article are supplied with analytic
proofs. No proof-assistant formalization is claimed. The basic Gamma
specialization is known and credited. The general harmonic-prefix
conjecture, $S_6$, and revised $S_8$ remain outside the proved scope
of this package. Formal coefficient calculations made during the
source search do not settle those period identities. The original
incoming archives and manuscript files are left unchanged.

\section{Further research questions}
\label{sec:questions}

\paragraph{1. Minimal compensation for higher resonant jets.}
The first-jet contact is a single Bernoulli constant. For
$\partial_s^d\Rcal_{N/2}(-N,w)$ with even $N$ and $d\ge2$,
logarithmic reciprocal tails appear. Determine a fixed symmetric
cutoff subtraction that evaluates these minimal jets, and derive
the complete contact polynomial in $\log\tau$ and the cutoff
normalization. Theorem~\ref{thm:negative-jets} supplies an ordinary
oversubtracted reference against which every proposed subtraction
can be tested.

\paragraph{2. Rationally related vertical scales.}
Replace $\zeta(r,a+\ii hn)\zeta(t,b-\ii hn)$ by
$\zeta(r,a+\ii h_1n)\zeta(t,b-\ii h_2n)$ with $h_1/h_2\in\mathbb Q$.
The correlation constraint is then oblique. Can each sampling image
still be reduced to a finite Hurwitz--Lerch grid, uniformly in both
integer orders? The continuous rational-scale results in
\cite{RepoVertical} suggest a precise comparison, but do not supply
the discrete image terms automatically.

\paragraph{3. Bilinear critical-order contacts.}
The difference-pairing proof uses $\Ree s,\Ree t>1/2$.
Study continuation to $s+t=1$ with the pole cancellation and the
lattice cutoff fixed. Determine whether a finite local functional
at the Mellin origin controls the discrepancy, as in
Theorem~\ref{thm:negative}, and identify its complete mixed jet.

\paragraph{4. Higher products and image geometry.}
For three or more vertical factors, Poisson sampling constrains a
sum of Mellin variables to an integer multiple of $\tau$.
Determine finite formulas for integer orders on each resulting
simplex, and identify when the result closes in multiple Lerch
functions or ordinary multiple polylogarithms. A useful first case
is a product of two trigammas against one reflected trigamma.

\paragraph{5. Order derivatives of nonzero images.}
Equation~\eqref{eq:all-diff-jets} is explicit but retains logarithmic
integrals. Determine which combinations of mixed spectral jets
reduce to finitely many Lerch order derivatives, and which naturally
require a genuinely larger function class. A failure of a numerical
basis search would not prove nonreducibility.

\paragraph{6. Periodic filters beyond the mean mode.}
The first resonant contact is killed exactly by a zero-mean periodic
filter. For higher minimal jets, determine which additional discrete
moments of the filter kill the logarithmic contacts. Relate those
conditions to root-of-unity polylogarithm coordinates without
asserting arithmetic independence of the resulting values.

\paragraph{7. Modular transport of the Stieltjes jet family.}
The $a=1$ Gamma specialization is related to the eta transformation
in \cite{ArdehaliRosengren}. Investigate whether selected spectral
order derivatives of \eqref{eq:master} admit a useful exact
$\tau\leftrightarrow4\pi^2/\tau$ transport, including every
nonmodular correction. This should begin with a functional identity,
not a guessed numerical relation among special values.

\paragraph{8. Formal verification of the analytic interchange.}
A proof-assistant development could isolate the periodization lemma,
normal convergence of compensated Hurwitz jets, and the
one-extra-subtraction contact argument. These analytic dependencies
are more consequential than a large list of special-value checks.
The elementary Bell compiler is a separate, comparatively small
formal algebra problem.

\section{Conclusion}

The discrete vertical problem admits a single organizing mechanism:
remove the finitely many Mellin-origin terms that obstruct ordinary
summation, periodize the remaining density, and extract spectral
coefficients only on a normal convergence domain. This yields exact
Stieltjes sums, finite polylogarithmic differences, and normalized
primitives. At the first negative-even boundary jet, restoring one
removed reciprocal term produces an explicit Bernoulli contact.
The same periodization argument applied to correlations gives the
finite-image polygamma and four-shift formulas.

The separation between absolute sums, symmetric limits, and
continued spectral expressions is part of each identity, not an
optional interpretive convention. It is also what makes the
formulas usable as independently testable additions to the ProveIt
manuscript rather than numerical conjectures.

\appendix
\section{Reproduction and source map}

The package is self-contained. Compile the article with
\begin{verbatim}
python build.py
\end{verbatim}
and replay the mathematics with
\begin{verbatim}
python -m pip install -r requirements.txt
python verification/run_checks.py
\end{verbatim}
The build script runs three LaTeX passes and checks for undefined
references and overfull boxes. The delivered PDF additionally
received a rendered layout review. A rebuild may differ in binary
metadata; its bytes are not covered by the original checksum.

The machine-readable source audit records the individual paths and
blob identifiers returned during inspection. In particular, the
continuous vertical report was read at commit
\src{da2792fc9cc8ef19c6a314be6c65e1437fbdc66b}; other reads used
\src{main} and returned their own blob identifiers. The paper's
source map must not be interpreted as a complete checkout or an
audit of unseen ZIP members.

\begin{thebibliography}{99}
\bibitem{DLMF25}
NIST Digital Library of Mathematical Functions, \S25.11,
\emph{Hurwitz Zeta Function}. Definitions, Euler--Maclaurin formulas,
special values, and sectorial asymptotics.
\url{https://dlmf.nist.gov/25.11}. Accessed 11 October 2026.

\bibitem{DLMF5}
NIST Digital Library of Mathematical Functions, \S\S5.9 and 5.11,
\emph{Gamma Function: Integral Representations and Asymptotic Expansions}.
\url{https://dlmf.nist.gov/5.9}; \url{https://dlmf.nist.gov/5.11}.
Accessed 11 October 2026.

\bibitem{DLMFPoly}
NIST Digital Library of Mathematical Functions, \S25.12,
\emph{Polylogarithms}. \url{https://dlmf.nist.gov/25.12}.
Accessed 11 October 2026.

\bibitem{DLMFLerch}
NIST Digital Library of Mathematical Functions, \S25.14,
\emph{Lerch's Transcendent}. \url{https://dlmf.nist.gov/25.14}.
Accessed 11 October 2026.

\bibitem{DLMFFourier}
NIST Digital Library of Mathematical Functions, \S1.8,
\emph{Fourier Series}. \url{https://dlmf.nist.gov/1.8}.
Accessed 11 October 2026.

\bibitem{Coffey}
M. W. Coffey, \emph{Parameterized summation relations for the Stieltjes
constants}, arXiv:1002.4684, 2010.
\url{https://arxiv.org/abs/1002.4684}.

\bibitem{ArdehaliRosengren}
A. Arabi Ardehali and H. Rosengren,
\emph{A New Product Formula for $(z;q)_\infty$, with Applications to
Asymptotics}, Constructive Approximation, published 19 September 2026.
DOI: 10.1007/s00365-026-09783-2.
\url{https://doi.org/10.1007/s00365-026-09783-2}.
In particular, Eqs.~(9), (13), and the discussion of symmetric products.

\bibitem{RepoMain}
ProveIt Contributors, \emph{Polylogarithms and their Arithmetic Bridges},
collective manuscript README, repository
\src{VladimirReshetnikov/ProveIt}, path
\src{Analysis/Polylogarithms/docs/manuscript/README.md}.
Inspected blob \src{078716656e8d73e98415b927cca1a1bfea059134}.
\url{https://github.com/VladimirReshetnikov/ProveIt/tree/main/Analysis/Polylogarithms/docs/manuscript}.

\bibitem{RepoVertical}
ProveIt research report, \emph{Vertical-Line Hurwitz Pairings and Exact
Stieltjes--Gamma Calculus}, 10 October 2026, README and stated scope.
Path under \src{Analysis/Polylogarithms/docs/reports/}:
\src{stieltjes-correlation-calculus/continuations/vertical-stieltjes-pairings/README.md}.
Read at commit \src{da2792fc9cc8ef19c6a314be6c65e1437fbdc66b};
blob \src{022b7ad4b25e28fd36322dd8ad36edfcf0d36285}.

\bibitem{RepoBernoulli}
ProveIt research report, \emph{Bernoulli Determinants, Cubic Tornheim Jets,
and Canonical Digamma Integrals}, 11 October 2026, README and
\src{sections/03b_general_hankel.tex}.
Path under \src{Analysis/Polylogarithms/docs/reports/}:
\src{stieltjes-harmonic-resonance/continuations/bernoulli-tornheim-digamma/}.
Inspected README blob \src{395bdbb9469f09109b3dad1d1a86822fe2339ac7}.
The already proved smaller-order-two reconstruction is not claimed here.
\end{thebibliography}
\end{document}
